% Fichier complété par tools/complete_missing_corrections.py.
% Source : DYN/DYN-04-TorseurDynamique/Cours/Cours.tex
% Statut : rédigé (5 question(s)).
\begin{corrigebox}[Corrigé rédigé]
\CorrigeQuestion{1}
$\vec\sigma_A(S/0)=\vec\sigma_G(S/0)+
                \overrightarrow{AG}\wedge m\vec V_G$, avec
                $\vec\sigma_G=\mathbf I_G\vec\Omega$.
\CorrigeQuestion{2}
$\vec\delta_A(S/0)=d\vec\sigma_A/dt|_0$ et
                $\vec\delta_A=\vec\delta_G+
                \overrightarrow{AG}\wedge m\vec\Gamma_G+
                \vec V(A/0)\wedge m\vec V_G$. Pour un point $A$ fixe, le dernier terme
                est nul.
\CorrigeQuestion{3}
$\{\mathcal C(S/0)\}_A=\{m\vec V_G;\vec\sigma_A\}_A$.
\CorrigeQuestion{4}
$\{\mathcal D(S/0)\}_A=\{m\vec\Gamma_G;\vec\delta_A\}_A$.
\CorrigeQuestion{5}
Dans une base principale au centre d'inertie,
                $\mathbf I_G=\operatorname{diag}(A,B,C)$. En un autre point, on applique
                le théorème de Huygens.
\end{corrigebox}
